\documentclass[conference]{IEEEtran}
\usepackage{cite}
\usepackage{amsmath,amssymb,amsfonts}
\usepackage{algorithmic}
\usepackage{graphicx}
\usepackage{textcomp}
\usepackage{xcolor}
\usepackage{booktabs}
\usepackage{multirow}
\usepackage{hyperref}
\usepackage{balance}
\usepackage{float}
\usepackage{subcaption}
\usepackage{algorithm}
\definecolor{pendingcolor}{RGB}{200, 100, 50}
\newcommand{\pending}[1]{\textcolor{pendingcolor}{\textit{[#1]}}}

\begin{document}

\title{Does Simulator Fidelity Buy Transfer? A Leakage-Controlled Diagnosis of Generative Market World Models for Reinforcement-Learning Trading}

\author{
\IEEEauthorblockN{Prit Kudale\IEEEauthorrefmark{1}, Naman Dwivedi\IEEEauthorrefmark{2}, Raj Dandekar\IEEEauthorrefmark{2}, Rajat Dandekar\IEEEauthorrefmark{2}, Sreedath Panat\IEEEauthorrefmark{2}}
\IEEEauthorblockA{\IEEEauthorrefmark{1}prit.kudale@gmail.com \qquad \IEEEauthorrefmark{2}Vizuara AI Labs, hello@vizuara.com}
}

\maketitle

\begin{abstract}
Model-based reinforcement learning trains agents inside a learned simulator, or ``dream,'' which is appealing in finance where live exploration is costly. Yet generative-market research optimizes \emph{stylized-fact fidelity} (fat tails, volatility clustering, leverage), reinforcement-learning trading research is model-free, and the two have never been connected. We ask the question at their intersection: does a market simulator's stylized-fact fidelity predict how well an agent trained inside it transfers to real data? We build a leakage-controlled testbed: known stochastic-differential-equation Oracle markets, a dense capacity-matched \emph{fidelity dial}, a no-lookahead trading environment, and a Transformer mixture-density world model. Training Proximal Policy Optimization agents purely in the dream and backtesting on held-out data, we find that across the full resolved range of simulator fidelity the transfer Sharpe ratio (mean return divided by its standard deviation) is statistically flat (regression slope $95\%$ confidence intervals span zero) and never beats Buy-and-Hold: stylized-fact fidelity is \emph{decoupled} from transfer. A known-process decomposition shows the gap is not a world-model-quality problem; a perfect Oracle dream transfers no better, so the world-model bias is approximately zero. A controlled predictability experiment isolates the cause: when the simulator captures \emph{reward-relevant} directional structure, transfer rises in lockstep (slope $+2.83$), whereas a stylized-fact-faithful but signal-blind GARCH dream stays flat ($+0.48$, not significant). On real equities, with the world model trained on eighty S\&P large-caps and evaluated on the held-out future of the full universe, the same pattern holds: the dream agent (mean Sharpe $-0.66$), model-free PPO ($-0.19$), and a momentum baseline ($+0.05$) all trail Buy-and-Hold ($+0.63$). Every number traces to a released ledger, and a Deflated Sharpe over all dial trials is insignificant.
\end{abstract}

\begin{IEEEkeywords}
world models, model-based reinforcement learning, algorithmic trading, sim-to-real transfer, stylized facts, mixture density networks, backtest overfitting
\end{IEEEkeywords}

\section{Introduction}

Reinforcement learning (RL) agents can be trained inside a learned generative model of an environment rather than the environment itself. This idea, popularized as World Models~\cite{ha2018} and the Dreamer line~\cite{hafner2020,hafner2023}, is attractive in finance, where every interaction with the real market is costly and irreversible. If a faithful simulator of returns could be learned from historical data, an agent could be trained cheaply inside the simulator and then deployed. The natural prerequisite, pursued by a large generative-market literature, is to make the simulator \emph{realistic}: to reproduce the stylized facts of asset returns such as fat tails, volatility clustering, and the leverage effect, the last being the negative correlation between a return and future volatility~\cite{cont2001,takahashi2019,kwonlee2024,kim2025}.

A separate literature applies RL directly to trading, with no learned simulator, and has repeatedly found that reported ``profits'' are fragile artifacts of look-ahead bias and backtest overfitting~\cite{finrlmeta,gort2022,gym4real,yan2025}. These two literatures have never been connected, and a basic question sits unanswered at their intersection. Generative-market work optimizes stylized-fact fidelity; RL work needs an environment that yields a transferable policy. \emph{Does the first property predict the second?} There is reason to doubt it: stylized facts describe the \emph{marginal} distribution of returns and their volatility, whereas trading profit depends on the \emph{conditional mean}, that is the predictable direction of the next move, which we call the reward-relevant structure of the market and which stylized facts do not constrain. TransDreamer~\cite{chen2022} hints at the same conclusion from the RL side, observing that policy quality in a learned world model tracks the model's \emph{reward-prediction} accuracy rather than its observation fidelity. This hypothesis has never been tested in markets, and whether it holds is not obvious: a sufficiently faithful generative model might be expected to capture predictable structure as a byproduct.

We answer the question with a leakage-controlled diagnostic rather than a new trading algorithm. Figure~\ref{fig:overview} summarizes the study. We construct known stochastic-differential-equation (SDE) Oracle markets whose true dynamics we control, a dense and capacity-matched \emph{fidelity dial} that varies a simulator's stylized-fact fidelity while holding architecture and volatility fixed, a no-lookahead trading environment as the only evaluation harness, and a causal Transformer mixture-density world model. We train Proximal Policy Optimization (PPO)~\cite{ppo} agents purely on dream rollouts and backtest them on held-out real data. Because the Oracle is known, we can decompose the dream-to-real gap into world-model bias, finite-data error, and intrinsic non-stationarity. Because the testbed is leakage-controlled, every number is trustworthy and traceable.

Our central finding is that stylized-fact fidelity is \emph{decoupled} from transfer, and that the cause is the absence of reward-relevant structure, not poor simulator quality. We make four contributions.

\begin{enumerate}
\item \textbf{A leakage-controlled dream-to-real testbed} for market world models: known-SDE Oracle markets, a dense capacity-matched fidelity dial, a no-lookahead trading environment, a Transformer mixture-density world model, and an append-only results ledger in which every reported number is reproducible and a Deflated Sharpe is computed over all trials.
\item \textbf{The headline finding}: across the resolved range of simulator fidelity (from $0.34$ to $0.90$ on the fidelity score), the transfer-versus-fidelity regression slope has a confidence interval spanning zero and the agent never beats Buy-and-Hold. Under a pre-registered necessary/sufficient/decoupled rule, the verdict is \emph{decoupled} for both a stationary and a non-stationary generator.
\item \textbf{A mechanism}. A known-SDE decomposition shows the gap is not a world-model-quality problem (world-model bias $\approx 0$). A controlled predictability experiment then shows that transfer rises only when the simulator captures \emph{reward-relevant} directional structure (Oracle slope $+2.83$), while a stylized-fact-faithful but signal-blind GARCH dream stays flat ($+0.48$, not significant). This confirms the TransDreamer hypothesis in markets.
\item \textbf{Honest external validity}. On real S\&P large-caps, evaluated on the held-out future of the full universe, the dream and model-free agents both lose to Buy-and-Hold after costs, as does a trivial momentum baseline, and we explain why: the bottleneck is the absence of reward-relevant alpha, not simulator quality.
\end{enumerate}

\begin{figure*}[t]
\centering
\includegraphics[width=0.92\textwidth]{figures/output/fig1_overview.png}
\caption{Study overview. We train a PPO agent entirely inside a learned market world model whose stylized-fact fidelity is varied along a dial, then test transfer on a held-out real backtest. The result, established across synthetic Oracle markets and real equities, is that stylized-fact fidelity is decoupled from transfer and that reward-relevant structure is what transfers.}
\label{fig:overview}
\end{figure*}

The remainder of the paper reviews related work, details the testbed and method, presents the results, and discusses limitations.

\section{Background and Related Work}

\subsection{World models and model-based RL}
World Models~\cite{ha2018} trained a controller entirely inside a mixture-density recurrent dream and transferred it to the real environment, already documenting that a dream-trained policy can exploit world-model error off-distribution. The Dreamer line~\cite{hafner2020,hafner2023} learns a recurrent latent world model and an actor-critic purely by latent imagination. TransDreamer~\cite{chen2022} replaces the recurrent backbone with a Transformer state-space model and explicitly ties policy quality to the world model's reward-prediction accuracy rather than its observation fidelity, an observation our work tests directly. A recent study of model-based sim-to-real transfer~\cite{sim2real2025} reports that synthetic data can degrade a real policy and proposes a quantitative gap metric. All of this work is evaluated on games and control, never on financial markets.

\subsection{Generative market models and stylized facts}
Cont~\cite{cont2001} catalogued the model-free stylized facts of asset returns and the estimators for them, which remain the standard fidelity yardstick. Generative-adversarial~\cite{takahashi2019,kwonlee2024} and diffusion~\cite{kim2025} models are judged by how well their samples reproduce these facts. These are \emph{generators}, not interactive environments, and none is used to train an agent inside it. The mixture-density head we adopt for the world model originates with Bishop~\cite{bishop1994} and has been applied to financial risk with recurrent backbones~\cite{herrig2025}.

\subsection{RL for trading and the rigor problem}
FinRL-Meta~\cite{finrlmeta} provides a benchmark of gym-style market environments and names survivorship bias, backtest overfitting, and information leakage as first-order threats. Gort et al.~\cite{gort2022} apply combinatorial cross-validation and a probability-of-backtest-overfitting test and find that all RL agents lose out-of-sample, with PPO merely losing least. Gym4ReaL~\cite{gym4real} reports that PPO and DQN cannot consistently beat Buy-and-Hold on real foreign-exchange data. The Deflated Sharpe ratio~\cite{deflated} and the backtesting protocol of Arnott et al.~\cite{arnott2019} formalize the multiple-comparisons problem, and real-time machine-learning studies~\cite{yan2025} show that reported Sharpe ratios fall substantially once look-ahead bias is removed. Our anti-bluff protocol instantiates these defenses; our contribution is to ask the fidelity-to-transfer question these works never pose.

\section{Methodology}

\subsection{Overview}
We measure whether a market simulator's stylized-fact fidelity predicts the dream-to-real transfer of an agent trained inside it. The apparatus, shown in Figure~\ref{fig:method}, has five parts: known-SDE Oracle markets (Section~\ref{sec:markets}), a Transformer mixture-density world model with a recurrent ablation (Section~\ref{sec:wm}), a fidelity dial that varies simulator fidelity in a controlled way (Section~\ref{sec:dial}), a no-lookahead trading environment and the agent (Section~\ref{sec:env}), and the fidelity-to-transfer measurement with a reward-relevance probe (Section~\ref{sec:measure}).

\begin{figure*}[t]
\centering
\includegraphics[width=0.95\textwidth]{figures/output/fig2_method.png}
\caption{Method. A causal Transformer with a mixture-density head learns the next-return distribution and is wrapped as a DreamEnv with the same Gym interface as the real backtest; a PPO agent is trained only on dream rollouts and evaluated on a no-lookahead trading environment. The lower band shows the timing: the agent observes returns through $t-1$, chooses a position, and earns $r_t$, which is unseen at the decision.}
\label{fig:method}
\end{figure*}

\subsection{The no-lookahead trading environment}\label{sec:env}
The only evaluation harness is a trading environment over a sequence of log-returns $r_t = \log(p_t/p_{t-1})$. At step $t$ the agent observes the window of returns available through $t-1$, chooses a position $a_t \in \{-1, 0, +1\}$ (short, flat, long), and earns the next return $r_t$, which is unseen at the decision. The reward applies a proportional transaction cost on position changes,
\begin{equation}
R_t = a_t \, r_t - \kappa \, |a_t - a_{t-1}|,
\label{eq:reward}
\end{equation}
with $\kappa$ the per-unit cost. This timing convention, illustrated in the lower band of Figure~\ref{fig:method}, removes the canonical same-bar look-ahead leak. The identical environment is used for dream training (over simulator-generated paths) and for the real backtest (over held-out data), so the comparison isolates the simulator.

\subsection{Known-SDE Oracle markets}\label{sec:markets}
We generate markets from processes whose true dynamics are known, so that the Oracle simulator is exact and data are unlimited. Three generators span the relevant regimes. Geometric Brownian motion is a light-tailed control. A GJR-GARCH process with Student-$t$ innovations produces fat tails, volatility clustering, and the leverage effect,
\begin{equation}
\sigma_t^2 = \omega + (\alpha + \gamma\,\mathbf{1}[r_{t-1}<0])\,\varepsilon_{t-1}^2 + \beta\,\sigma_{t-1}^2,
\label{eq:garch}
\end{equation}
with $r_t = \mu + \sigma_t \epsilon_t$ and $\epsilon_t$ standardized Student-$t$. A two-state Markov-switching variant of \eqref{eq:garch} adds non-stationarity and is strictly richer than any single-regime GARCH. Figure~\ref{fig:facts} confirms that the generators span light- to heavy-tailed returns and bracket the stylized facts of real equities.

\subsection{Transformer mixture-density world model}\label{sec:wm}
The world model predicts the conditional distribution of the next return from a length-$L$ context window. A causal Transformer encoder with masked self-attention maps the window to a hidden state, and a mixture-density head emits a $K$-component Gaussian mixture with softmax weights $\pi_k$, linear means $\mu_k$, and softplus-style positive scales $\sigma_k$. The model is trained by minimizing the mixture negative log-likelihood with a small penalty that discourages component collapse,
\begin{equation}
\mathcal{L} = -\,\mathbb{E}\!\left[\log \sum_{k=1}^{K} \pi_k \, \mathcal{N}(r;\mu_k,\sigma_k^2)\right] + \lambda \sum_k \pi_k^2 .
\label{eq:mdn}
\end{equation}
A long short-term memory backbone with the same head is the recurrent ablation. The world model is fit on the training split only, and dream paths are generated autoregressively by sampling from \eqref{eq:mdn} starting from training-split context windows.

\subsection{The fidelity dial}\label{sec:dial}
To make the fidelity-to-transfer question falsifiable rather than a comparison of a few confounded models, we vary fidelity continuously at a fixed parametric class. A scalar $\eta \in [0,1]$ interpolates a simulator's parameters from a Gaussian-degenerate setting ($\eta=0$, no tails, clustering, or leverage) to the Oracle-matched setting ($\eta=1$),
\begin{equation}
\theta(\eta) = (1-\eta)\,\theta_{\text{Gauss}} + \eta\,\theta_{\text{Oracle}},
\label{eq:dial}
\end{equation}
after which the output is rescaled to the Oracle's unconditional standard deviation. Because the four discriminating stylized facts are scale-invariant, this rescaling removes any volatility-level confound: along the dial only the \emph{shape} of returns changes, not their magnitude.

\subsection{Measuring fidelity and transfer}\label{sec:measure}
Simulator fidelity is a bounded score over the four facts,
\begin{equation}
F = \exp\!\Big(-\tfrac{1}{4}\textstyle\sum_{i} \min(|s_i^{\text{sim}} - s_i^{\text{real}}|/c_i,\,3)\Big),
\label{eq:fidelity}
\end{equation}
where $s_i$ are excess kurtosis, Hill tail index, the autocorrelation of absolute returns, and the leverage correlation, and $c_i$ are fact-specific scales. For each dial setting we wrap the simulator as a DreamEnv, train PPO only on dream rollouts, and record the realized transfer Sharpe on the held-out test path. We regress transfer on $F$ across the dial and apply a pre-registered rule: the relationship is \emph{necessary} if the slope confidence interval excludes zero and the low-fidelity end matches Random, \emph{sufficient} if high-fidelity transfer beats Buy-and-Hold and is non-inferior to model-free PPO, and \emph{decoupled} if the slope interval includes zero over a resolved fidelity range and high fidelity does not beat Buy-and-Hold. A Deflated Sharpe ratio over all trials guards against multiple comparisons.

To probe the mechanism we add a market with a tunable, learnable directional edge. A hidden persistent trend $s_t$ is added to the conditional mean on top of the GJR-GARCH volatility,
\begin{equation}
s_t = \psi\,s_{t-1} + \sigma_s\,\xi_t, \qquad r_t = \mu + s_t + \sigma_t\,\epsilon_t,
\label{eq:momentum}
\end{equation}
with persistence $\psi$ close to one, so that $\mathbb{E}[r_t \mid \text{history}] \approx s_t$ is inferable from recent returns; the trend variance $\sigma_s^2$ sets the signal-to-noise ratio and hence the realizable directional $R^2$. On this market we compare three dreams trained with identical budgets: the Oracle (which contains the trend), a GARCH fit (which reproduces the stylized facts but has a constant conditional mean, hence is signal-blind), and a Gaussian floor.

\section{Experimental Setup}

\subsection{Data}
The synthetic arm uses the three generators with strictly chronological $60{,}000 / 20{,}000 / 20{,}000$ train/validation/test splits per seed. The real arm pools eighty liquid large-cap United States equities (S\&P~500 constituents) over 2010--2024; point-in-time index membership is not controlled, so real-arm results are an optimistic bound. The real splits are chronological and cross-sectional, holding out both three tickers and the final time period; the shared-backbone world model never sees the held-out tickers or the held-out future.

\subsection{Implementation}
Experiments run on Modal serverless GPUs, with L4 and A10G for the synthetic arm and an A100-80GB for the pooled real backbone. The world model is implemented in PyTorch and the agent uses Proximal Policy Optimization from Stable-Baselines3~\cite{sb3}; GARCH baselines use the \texttt{arch} package. Leakage is enforced in code: normalization statistics are computed on the training split only, dream rollouts are seeded only from training windows, and the world model never sees test data. PPO hyperparameters are identical across all dream, model-free, and Oracle conditions; the per-experiment step budget is matched within every head-to-head comparison (200{,}000 steps for the dial and predictability experiments, 150{,}000 for the learning-curve and ablation experiments). Hyperparameters are listed in Table~\ref{tab:hyper}. All comparisons use at least three seeds with bootstrap confidence intervals. The code, the world model, and the append-only results ledger from which every reported number is reproducible are publicly released.\footnote{\url{https://github.com/OmuNaman/dream-to-real-trading}}

\begin{table}[!t]
\caption{Hyperparameters and synthetic-generator parameters. All values are fixed across conditions; the per-experiment PPO step budget is matched within every head-to-head comparison.}
\label{tab:hyper}
\centering
\resizebox{\columnwidth}{!}{%
\begin{tabular}{@{}lll@{}}
\toprule
Component & Setting & Value \\
\midrule
World model & layers / width / mixture $K$ & 2 / 64 / 5 \\
            & context $L$ / MDN penalty $\lambda$ & 60 / 0.1 \\
            & optimizer / LR / epochs & Adam / $10^{-3}$ / 20 \\
PPO         & policy / value net & MLP $[64,64]$ \\
            & rollout / batch / epochs & 2048 / 256 / 10 \\
            & $\gamma$ / GAE $\lambda$ / clip & 0.99 / 0.95 / 0.2 \\
            & learning rate / entropy & $3{\times}10^{-4}$ / 0.01 \\
            & steps (dial / learning) & 200k / 150k \\
Environment & actions / cost $\kappa$ & \{short,flat,long\} / 5 bps \\
Generators  & GJR $(\omega,\alpha,\gamma,\beta,\nu)$ & $(2{\times}10^{-6},.04,.12,.88,5)$ \\
            & regime $\beta_{\text{lo}},\beta_{\text{hi}}$, vol mult. & $0.83,\,0.88,\,4{\times}$ \\
            & trend $\psi$ / momentum $L$ & 0.96 / 60 \\
\bottomrule
\end{tabular}}
\end{table}

\subsection{Baselines}
Buy-and-Hold and a Random policy are the passive references. Model-free PPO, trained directly on the real training series with the same environment, budget, and seeds, is the comparator that isolates the value of the simulator. The fidelity ladder additionally includes a Gaussian floor, a GJR-GARCH simulator, and, in the synthetic arm, the Oracle world model as a perfect-fidelity ceiling.

\begin{figure}[t]
\centering
\includegraphics[width=\columnwidth]{figures/output/fig3_stylized_facts.png}
\caption{The synthetic generators span light- to heavy-tailed returns and reproduce the stylized facts of real equities, with synthetic tails heavier than the real value rather than centered on it (excess kurtosis on a log scale, tail index, volatility clustering, and the leverage correlation).}
\label{fig:facts}
\end{figure}

\section{Results}

\subsection{World-model fidelity}
Table~\ref{tab:fidelity} reports held-out fidelity by simulator and generator. On the light-tailed control all simulators tie. On the GJR-GARCH process the neural simulators recover the true class and clear the Gaussian floor (fidelity $0.34$ versus $0.63$--$0.70$). On the non-stationary regime-switching process the neural mixture-density models reproduce the facts well (Transformer and recurrent both near $0.72$, range $0.66$--$0.79$) while a single GARCH collapses to $0.23$--$0.30$, because fitting a stationary GARCH to regime-switching data drives it toward an integrated regime whose simulated kurtosis explodes. The held-out negative log-likelihood, reported in raw return units so that the learned mixture is compared fairly against the Gaussian and GARCH baselines, ranks the Gaussian floor ($-3.26$) well below the neural and GARCH models (near $-3.46$). Figure~\ref{fig:facts} shows that the generators reproduce the same qualitative facts as real equities, with synthetic tails heavier than the real value rather than centered on it.

\begin{table}[t]
\caption{Held-out stylized-fact fidelity $F$ (\eqref{eq:fidelity}, $\uparrow$) by simulator and generator, best per generator in bold. On its own class (GJR-GARCH) a single GARCH wins, but it collapses on non-stationary (regime-switch) data, where the Transformer mixture-density model matches the recurrent one and is the best robust simulator.}
\label{tab:fidelity}
\centering
\small
\begin{tabular}{@{}lccc@{}}
\toprule
Simulator & GBM & GJR-GARCH & Regime-switch \\
\midrule
Gaussian floor   & 0.96 & 0.34 & 0.34 \\
GARCH            & 0.96 & \textbf{0.70} & 0.27 \\
RNN-MDN          & 0.90 & 0.63 & 0.72 \\
Transformer-MDN  & 0.93 & 0.65 & \textbf{0.72} \\
\bottomrule
\end{tabular}
\end{table}

\subsection{Fidelity is decoupled from transfer}
Figure~\ref{fig:decoupled} is the central result. Across the dense dial, the regression of realized transfer Sharpe on simulator fidelity is statistically flat: the slope is $+0.20$ with a $95\%$ confidence interval of $[-0.37, 0.60]$ on GJR-GARCH and $-0.09$ with interval $[-0.47, 0.23]$ on the regime-switching process, both spanning zero over resolved fidelity ranges of $0.39$ and $0.54$ (the dial moves fidelity from $0.34$ to $0.73$ and from $0.35$ to $0.90$ respectively). The dream-trained agent earns a transfer Sharpe near $-0.5$ at every fidelity level, above the Random policy ($-0.96$ and $-0.72$) but always below Buy-and-Hold ($+0.41$ and $+0.33$). Under the pre-registered rule the verdict is \emph{decoupled} for both generators: higher stylized-fact fidelity does not buy better transfer. The mean in-dream-to-real gap is small (typically within $\pm 0.15$ Sharpe), so the agent is not merely overfitting the dream; it has no edge to overfit. Over the 2010--2024 sample Buy-and-Hold is a strong long-only benchmark on a trending equity market, so the regime-independent finding is the \emph{decoupling} itself, the flat slope whose confidence interval spans zero, rather than the comparison to Buy-and-Hold.

\begin{figure*}[t]
\centering
\includegraphics[width=0.95\textwidth]{figures/output/fig4_decoupled.png}
\caption{Headline. Realized dream-to-real transfer Sharpe regressed on simulator stylized-fact fidelity for the stationary (left) and non-stationary (right) generator, over resolved fidelity ranges of $0.39$ and $0.54$. Both slopes have confidence intervals spanning zero, and the dream agent never beats Buy-and-Hold. Stylized-fact fidelity is decoupled from transfer.}
\label{fig:decoupled}
\end{figure*}

\subsection{Reward-relevant structure is what transfers}
If simulator quality is not the bottleneck, what is? Figure~\ref{fig:mechanism} answers. On a market with a tunable directional edge, the Oracle dream, which captures the trend, transfers in lockstep with predictability: its transfer Sharpe rises monotonically from $-0.66$ at zero edge to $+13.1$ at the maximum realizable directional $R^2$ of $0.22$. Regressed against the signal-to-noise dial (Fig.~\ref{fig:mechanism}), the slope is $+2.83$ with confidence interval $[2.16, 4.99]$. A GARCH dream that reproduces the same stylized facts (signal-averaged fidelity $0.29$--$0.74$) but has a constant conditional mean is signal-blind and stays flat, with a slope of $+0.48$ and interval $[-0.28, 1.49]$ that includes zero. The Gaussian floor and Buy-and-Hold are likewise flat. The large Oracle Sharpe values are a controlled synthetic edge by construction, not a tradeable claim; the result of record is the contrast between a rising Oracle and a flat GARCH. This experiment also demonstrates that the agent is not broken: the same PPO pipeline that produces flat transfer on stylized-fact-only markets finds a strong edge whenever a reward-relevant one exists, which rules out a weak policy class as the explanation for the flat transfer. Reward-relevant conditional accuracy, not marginal stylized-fact fidelity, is what buys transfer.

\begin{figure}[t]
\centering
\includegraphics[width=\columnwidth]{figures/output/fig5_mechanism.png}
\caption{The mechanism. When the dream captures reward-relevant directional structure (Oracle), transfer rises with predictability (slope $+2.83$, CI $[2.16,4.99]$); when it reproduces stylized facts but is signal-blind (GARCH), transfer stays flat (slope $+0.48$, CI $[-0.28,1.49]$, not significant) despite high fidelity. The large Oracle Sharpe values are a controlled synthetic edge by construction and are not a tradeable result; the contrast between the rising Oracle and the flat GARCH is the finding.}
\label{fig:mechanism}
\end{figure}

\subsection{The gap is not a world-model-quality problem}
A natural explanation for flat transfer would be a poor world model, but the known-SDE decomposition rules this out. Figure~\ref{fig:decomp}(a) shows that a perfect Oracle dream ($-0.67$ on GJR-GARCH) transfers no better than the learned dream ($-0.48$) or model-free PPO ($-0.64$), and none beats Buy-and-Hold; the world-model bias, defined as Oracle-dream minus learned-dream transfer, is only $-0.19$ on GJR-GARCH and $-0.04$ on the regime process. More training data raises fidelity but barely moves transfer (a finite-data effect of $+0.23$ Sharpe across a thirty-fold increase in data, still negative), and the non-stationarity effect is negligible ($-0.05$). None of world-model bias, finite data, or non-stationarity explains the gap. As a separate backbone-efficiency observation, Figure~\ref{fig:decomp}(b) shows that the recurrent and Transformer mixture-density models cross over with data: the recurrent model is more fidelity-efficient at small data ($0.556$ versus $0.503$ at 2{,}000 bars) while the Transformer overtakes it at 60{,}000 bars ($0.661$ versus $0.631$).

\begin{figure*}[t]
\centering
\includegraphics[width=0.95\textwidth]{figures/output/fig6_decomposition.png}
\caption{(a) Decomposition of the dream-to-real gap: a perfect Oracle dream transfers no better than the learned dream or model-free PPO, and all lose to Buy-and-Hold, so the world-model bias is near zero. (b) A separate backbone-efficiency result: the recurrent and Transformer mixture-density models cross over with data, the Transformer overtaking at sixty thousand bars.}
\label{fig:decomp}
\end{figure*}

\subsection{Costs, significance, and real-market validity}
Figure~\ref{fig:costreal}(a) confirms that the negative synthetic transfer is transaction-cost bleed rather than anti-skill: at zero cost the high-fidelity dream is approximately flat ($-0.03$ Sharpe), falling to $-0.73$ at 5 basis points and $-1.64$ at 20 basis points. With no directional alpha to extract, any trading simply pays the cost. A Deflated Sharpe ratio over all sixty-nine dial trials, fed the true trial count, gives a probability of $1.4\times 10^{-9}$ that the best Sharpe ($0.044$) is genuine, so no agent has a real edge after correcting for multiple comparisons. On real equities (Figure~\ref{fig:costreal}(b)) the Transformer world model, trained on eighty S\&P large-caps, reproduces real stylized facts (real excess kurtosis $23.7$, tail index $2.67$, clustering $0.25$, leverage $-0.04$) at fidelity $0.63$. Evaluated on the held-out future of the full universe (eighty in-sample tickers tested out-of-time plus three fully held-out tickers, $83$ series in total), the dream agent (mean Sharpe $-0.66$) and model-free PPO ($-0.19$) both lose to Buy-and-Hold ($+0.63$), as does a trivial time-series momentum baseline ($+0.05$), confirming that the failure is not specific to the reinforcement-learning policy class. A date-aligned panel bootstrap, with an effective sample size of $16.7$ after accounting for cross-sectional correlation among the strategy series, gives a pooled equal-weight dream Sharpe of $-2.42$ with interval $[-3.64, -1.06]$, significantly negative. The synthetic decoupling and the no-edge failure both hold on real data.

\begin{figure*}[t]
\centering
\includegraphics[width=0.95\textwidth]{figures/output/fig7_cost_real.png}
\caption{(a) The high-fidelity dream is approximately flat at zero transaction cost; realistic costs turn trading without an edge into pure bleed. (b) On real S\&P large-caps, the dream, model-free, and momentum policies all lose to Buy-and-Hold; a representative subset of tickers is shown (ten of the $83$ series; held-out tickers starred), while the reported headline numbers are pooled over the held-out future of the full universe.}
\label{fig:costreal}
\end{figure*}

\subsection{Discussion}
The results have a direct implication for using generative-market models as RL training environments. Stylized-fact fidelity, the property that generative-market research optimizes, measures the realism of the return \emph{marginal} and its volatility structure. Trading profit depends on the \emph{conditional mean}, which stylized facts do not constrain. A simulator can therefore be highly faithful and still useless for training a profitable trader, exactly as the signal-blind GARCH dream demonstrates. The right objective for a market world model intended for control is reward-relevant conditional accuracy, not marginal fidelity. The dream-to-real gap in trading is a task-structure problem, not a simulator-quality problem.

\section{Limitations and Future Work}
The study is deliberately scoped. The action space is single-asset and discrete, and the synthetic markets are daily-frequency; intraday and limit-order-book dynamics are left to future work. The fidelity score is a composite over four facts, and the synthetic edge in the predictability experiment is injected by construction rather than discovered, so its large Sharpe values are illustrative and not tradeable. The real arm uses a current liquid universe without point-in-time membership, so its numbers are an optimistic bound, and a single seed supported by a cross-sectional panel bootstrap. The held-out world-model likelihoods use best-validation early stopping and are mildly optimistic, though identically for both backbones. The natural next step is to optimize world models for reward-relevant conditional accuracy directly and to test whether that, unlike stylized-fact fidelity, buys transfer.

\section{Conclusion}
Stylized-fact fidelity, the property generative-market models are built to optimize, does not buy dream-to-real transfer for reinforcement-learning traders. Across a controlled range of simulator fidelity, transfer is flat and never beats Buy-and-Hold, and a known-SDE decomposition shows the gap is not a world-model-quality problem. Reward-relevant conditional accuracy is what transfers: when a simulator captures directional structure, transfer rises in lockstep, while a stylized-fact-faithful but signal-blind simulator does not transfer at all. On real equities the same pattern holds. The dream-to-real gap in trading is a task-structure problem, and we characterize it cleanly under leakage control with every number traceable to a released ledger.

\section*{Acknowledgments}
The authors thank Vizuara AI Labs for providing mentorship and computational resources for this research.

\balance
\begin{thebibliography}{99}
\bibitem{ha2018} D. Ha and J. Schmidhuber, ``Recurrent world models facilitate policy evolution,'' \textit{Advances in Neural Information Processing Systems (NeurIPS)}, vol. 31, 2018.
\bibitem{hafner2020} D. Hafner, T. Lillicrap, J. Ba, and M. Norouzi, ``Dream to control: Learning behaviors by latent imagination,'' \textit{International Conference on Learning Representations (ICLR)}, 2020.
\bibitem{hafner2023} D. Hafner, J. Pasukonis, J. Ba, and T. Lillicrap, ``Mastering diverse domains through world models,'' \textit{arXiv:2301.04104}, 2023.
\bibitem{chen2022} C. Chen, Y.-F. Wu, J. Yoon, and S. Ahn, ``TransDreamer: Reinforcement learning with Transformer world models,'' \textit{arXiv:2202.09481}, 2022.
\bibitem{sim2real2025} Z. Lin and S. Sun, ``Revealing the challenges of sim-to-real transfer in model-based reinforcement learning via latent space modeling,'' \textit{arXiv:2506.12735}, 2025.
\bibitem{cont2001} R. Cont, ``Empirical properties of asset returns: Stylized facts and statistical issues,'' \textit{Quantitative Finance}, vol. 1, no. 2, pp. 223--236, 2001.
\bibitem{takahashi2019} S. Takahashi, Y. Chen, and K. Tanaka-Ishii, ``Modeling financial time-series with generative adversarial networks,'' \textit{Physica A}, vol. 527, 121261, 2019.
\bibitem{kwonlee2024} S. Kwon and Y. Lee, ``Can GANs learn the stylized facts of financial time series?,'' \textit{ACM Int. Conf. on AI in Finance (ICAIF)}, 2024.
\bibitem{kim2025} G. Kim, S.-Y. Choi, and Y. Kim, ``A diffusion-based generative model for financial time series via geometric Brownian motion,'' \textit{arXiv:2507.19003}, 2025.
\bibitem{bishop1994} C. M. Bishop, ``Mixture density networks,'' \textit{Technical Report NCRG/94/004}, Aston University, 1994.
\bibitem{herrig2025} N. Herrig, ``Risk forecasting using long short-term memory mixture density networks,'' \textit{arXiv:2501.01278}, 2025.
\bibitem{finrlmeta} X.-Y. Liu et al., ``FinRL-Meta: Market environments and benchmarks for data-driven financial reinforcement learning,'' \textit{NeurIPS Datasets and Benchmarks}, 2022.
\bibitem{gort2022} B. J. D. Gort et al., ``Deep reinforcement learning for cryptocurrency trading: Practical approach to address backtest overfitting,'' \textit{ICAIF Workshop}, 2022.
\bibitem{gym4real} D. Salaorni et al., ``Gym4ReaL: A suite for benchmarking real-world reinforcement learning,'' \textit{arXiv:2507.00257}, 2025.
\bibitem{ppo} J. Schulman, F. Wolski, P. Dhariwal, A. Radford, and O. Klimov, ``Proximal policy optimization algorithms,'' \textit{arXiv:1707.06347}, 2017.
\bibitem{sb3} A. Raffin et al., ``Stable-Baselines3: Reliable reinforcement learning implementations,'' \textit{Journal of Machine Learning Research}, vol. 22, no. 268, pp. 1--8, 2021.
\bibitem{deflated} D. H. Bailey and M. L\'opez de Prado, ``The deflated Sharpe ratio: Correcting for selection bias, backtest overfitting, and non-normality,'' \textit{Journal of Portfolio Management}, vol. 40, no. 5, pp. 94--107, 2014.
\bibitem{arnott2019} R. Arnott, C. R. Harvey, and H. Markowitz, ``A backtesting protocol in the era of machine learning,'' \textit{Journal of Financial Data Science}, vol. 1, no. 1, pp. 64--74, 2019.
\bibitem{yan2025} X. Yan, B. Li, A. G. Rossi, and L. Zheng, ``Machine learning from a universe of signals: The role of feature engineering,'' \textit{Journal of Financial Economics}, 2025.
\bibitem{gukellyxiu2020} S. Gu, B. Kelly, and D. Xiu, ``Empirical asset pricing via machine learning,'' \textit{Review of Financial Studies}, vol. 33, no. 5, pp. 2223--2273, 2020.
\end{thebibliography}

\end{document}
